\section{Minor arcs}\label{sec:minor}

Use canonical representatives $0\le h<L$.  Since $2N+1\le L$, the residues of
the integers $|m|\le N$ are distinct.  Let $S_M$ be this major-arc set and put
\[
 S_m=(\Z/L\Z)\setminus S_M.
\]
We divide the minor arcs into two disjoint classes:
\begin{equation}\label{eq:minor-lanes}
 S_{\mathrm{blk}}
 =S_m\cap\{h:a(h)\notin\mathfrak M(C)\},
 \qquad
 S_{\mathrm{ext}}
 =S_m\cap\{h:a(h)\in\mathfrak M(C)\}.
\end{equation}
Thus $S_{\mathrm{blk}}\dot\cup S_{\mathrm{ext}}=S_m$.

Put
\[
 P=\prod_{p\in\Bset}p,
 \qquad L=bP.
\]
If two frequencies have the same block assignment, their difference is
divisible by $P$.  In $[0,bP)$ a fixed residue modulo $P$ has at most the $b$
representatives $r,r+P,\ldots,r+(b-1)P$.  Since
$\Qctrl(a(h))\le\QE(h)$, Theorem~\ref{thm:global-partition} with $c=16/9$
gives
\begin{equation}\label{eq:block-minor-budget}
 \sum_{h\in S_{\mathrm{blk}}}|F(h)|
 \le
 \frac{b\bigl(\eta+C_{\mathrm{tail}}
 e^{-(16/9)C^2/2}\bigr)}{\sctrl}.
\end{equation}
The factor $b$ is the fibre multiplicity.

Now let $h\in S_{\mathrm{ext}}$.  There is a label $m=m(h)$ with
\[
 |m|\le C/\sctrl\le N,
 \qquad h\equiv m\pmod p\quad(p\in\Bset).
\]
Because $h\notin S_M$, it is not congruent to $m$ modulo $L$.  If it agreed
with $m$ modulo every prime divisor of the squarefree integer $b$, then it
would agree modulo $b$.  Together with the block congruences and the
coprimality of $b$ and $P$, the Chinese remainder theorem would give agreement
modulo $L$, a contradiction.  Hence there is a prime
\begin{equation}\label{eq:mismatch-prime}
 r(h)\mid b,
 \qquad h\not\equiv m(h)\pmod{r(h)}.
\end{equation}
This is the second use of the squarefreeness of $b$.

Let $G$ be the integer chosen in Section~\ref{sec:completion}.  For sufficiently
large $k_0$, select a set $\mathcal G$ of exactly $G$ primes in a high block,
all outside the prime support of $b$, with
\begin{equation}\label{eq:high-prime-separation}
 2N<s\qquad(s\in\mathcal G).
\end{equation}
The density estimate in Definition~\ref{def:AI} supplies these primes for sufficiently large $k_0$.  Let $\Eaux$ contain every auxiliary denominator
\begin{equation}\label{eq:aux-family}
 rs\qquad(r\mid b\text{ prime},\ s\in\mathcal G).
\end{equation}
Unique factorization recovers both primes, so these products are distinct; they
cannot coincide with products of two block primes because the block primes lie
outside the prime support of $b$.

Fix $h\in S_{\mathrm{ext}}$, write $m=m(h)$ and $r=r(h)$, and take
$s\in\mathcal G$.  Since $h\equiv m\pmod s$, write $h=m+js$.  The mismatch
modulo $r$ gives $j\not\equiv0\pmod r$.  Therefore
\[
 \normdist{j/r}\ge\frac1r,
 \qquad
 \left|\frac{m}{rs}\right|<\frac1{2r},
\]
and the triangle inequality on $\mathbb R/\mathbb Z$ gives
\[
 \normdist{h/(rs)}\ge\frac1{2r}.
\]
The one-factor identity used above now yields
\begin{equation}\label{eq:aux-factor}
 \left|(1-\theta)+\theta e^{2\pi ih/(rs)}\right|
 \le\sqrt{1-\frac8{9r^2}}
 \le\sqrt{1-\frac8{9b^2}}=: \beta_b<1.
\end{equation}
The $G$ factors are distinct, so $|F(h)|\le\beta_b^G$.  There are at most
$2N+1$ labels and at most $b$ frequencies in each block fibre; hence
\begin{equation}\label{eq:extra-count}
 |S_{\mathrm{ext}}|\le b(2N+1),
 \qquad
 \sum_{h\in S_{\mathrm{ext}}}|F(h)|
 \le b(2N+1)\beta_b^G.
\end{equation}
